\documentclass[11pt,a4paper]{article}
\usepackage[spanish,es-noquoting]{babel}
\usepackage[top=2.2cm, bottom=2.2cm, left=2.2cm, right=2.2cm]{geometry}
\usepackage{fontspec}
\usepackage{amsmath,amssymb}
\usepackage{enumitem}
\usepackage{microtype}

\setmainfont{Calibri}
\linespread{1.15}

\begin{document}
\pagestyle{empty}

\noindent\textbf{1.} Se desea fabricar un depósito cilíndrico cerrado con una capacidad de $16\pi\text{ m}^3$. El coste del material utilizado para la tapa superior es de $3\text{ €/m}^2$, el del suelo o base inferior es de $5\text{ €/m}^2$, y el de la pared lateral es de $4\text{ €/m}^2$.
\begin{enumerate}[label=\textbf{\alph*)}, leftmargin=2em, itemsep=6pt]
    \item Expresa la función del coste total de fabricación, $C(r)$, en función únicamente del radio $r$ de la base.
    \item Determina las dimensiones que debe tener el depósito (radio $r$ y altura $h$) para que el coste total sea mínimo.
    \item Justifica, mediante el criterio de la segunda derivada, que se trata efectivamente de un mínimo de coste y calcula el valor de dicho coste mínimo.
\end{enumerate}

\vspace{1cm}

\noindent\textbf{2.} Se desea fabricar un depósito cilíndrico cerrado con una capacidad de $54\pi\text{ m}^3$. El coste del material utilizado para la tapa superior es de $4\text{ €/m}^2$, el del suelo o base inferior es de $8\text{ €/m}^2$, y el de la pared lateral es de $6\text{ €/m}^2$.
\begin{enumerate}[label=\textbf{\alph*)}, leftmargin=2em, itemsep=6pt]
    \item Expresa la función del coste total de fabricación, $C(r)$, en función únicamente del radio $r$ de la base.
    \item Determina las dimensiones que debe tener el depósito (radio $r$ y altura $h$) para que el coste total sea mínimo.
    \item Justifica, mediante el criterio de la segunda derivada, que se trata efectivamente de un mínimo de coste y calcula el valor de dicho coste mínimo.
\end{enumerate}

\end{document}
